\subsection{激光超声NDT基础知识}
\label{sec:lu-fundamentals}

激光超声检测通过材料表面的快速热膨胀产生弹性波。两种产生机制共存：热弹性消融，其中激光脉冲引起应力而无材料移除；
以及烧蚀产生，其中等离子体形成喷射表面材料以产生更强的声波\cite{liu2022deep}。由此产生的波形是宽带的，包含频散Lamb波模态，其相速度随频率和板厚变化。这种波动特性使得薄壁结构中的缺陷检测成为可能，但模态叠加使信号解释变得复杂\cite{zhang2020deep}。

常见检测目标包括铝合金和碳纤维增强聚合物(CFRP)复合材料，激光超声可检测其中的分层、空洞和粘接线缺陷，无需耦合剂。
然而，由于多种噪声源，获取的信号存在低信噪比(SNR)：激光感应等离子体发射、环境声学干扰和光电探测器噪声。
因此，信号平均是标准做法——重复采集进行集合平均以抑制随机噪声并提高信噪比\cite{liu2022deep}。这种平均方法线性增加检测时间，并可能对敏感材料造成表面改性，激励了对单次去噪替代方案的研究。
